\section{Orden topol\'ogico}
\whenuse{Procesar dependencias de un DAG, hacer DP en grafo dirigido o detectar
ciclos comprobando si se ordenaron todos los v\'ertices.}

Kahn devuelve un orden solamente si el grafo dirigido es un DAG.

\begin{lstlisting}[style=cpp]
vector<int> topological_sort(const vector<vector<int>> &graph) {
    int n = graph.size();
    vector<int> indegree(n, 0), order;
    for (int v = 0; v < n; ++v)
        for (int to : graph[v]) ++indegree[to];

    queue<int> q;
    for (int v = 0; v < n; ++v) if (indegree[v] == 0) q.push(v);
    while (!q.empty()) {
        int v = q.front(); q.pop();
        order.push_back(v);
        for (int to : graph[v]) if (--indegree[to] == 0) q.push(to);
    }
    if ((int)order.size() != n) return {}; // existe un ciclo
    return order;
}
\end{lstlisting}

\complexity{$O(V+E)$ tiempo y memoria.}

\section{Componentes fuertemente conexas (Kosaraju)}
\whenuse{Contraer ciclos de un grafo dirigido, obtener su DAG de componentes o
resolver problemas de alcanzabilidad y 2-SAT.}

\begin{lstlisting}[style=cpp]
vector<int> strongly_connected_components(
        const vector<vector<int>> &graph) {
    int n = graph.size();
    vector<vector<int>> reversed(n);
    for (int v = 0; v < n; ++v)
        for (int to : graph[v]) reversed[to].push_back(v);

    vector<char> visited(n, false);
    vector<int> order;
    function<void(int)> dfs1 = [&](int v) {
        visited[v] = true;
        for (int to : graph[v]) if (!visited[to]) dfs1(to);
        order.push_back(v);
    };
    for (int v = 0; v < n; ++v) if (!visited[v]) dfs1(v);

    vector<int> component(n, -1);
    function<void(int,int)> dfs2 = [&](int v, int id) {
        component[v] = id;
        for (int to : reversed[v]) if (component[to] == -1) dfs2(to, id);
    };
    reverse(order.begin(), order.end());
    int count = 0;
    for (int v : order) if (component[v] == -1) dfs2(v, count++);
    return component;
}
\end{lstlisting}

\complexity{$O(V+E)$ tiempo y memoria.}
\warning{Las DFS recursivas pueden desbordar el stack para grafos muy grandes.}
